Diseño y prueba de un amperímetro analógico de tres rangos (10, 50 y 100 mA) construido a partir de un galvanómetro. Se cumplieron los objetivos de calibrar el instrumento y verificar sus principios de funcionamiento. La calibración se fijó en el rango de 100 mA, lo que permitió medir el error sistemático inducido en los rangos inferiores, resultando en desviaciones significativas de $\approx 14\%$ en los fondos de escala de 10 mA y 50 mA. Se comprobó exitosamente la linealidad del instrumento, con un $R^2 > 0.996$ en todos los rangos, validando el diseño teórico. Finalmente, se cuantificó el efecto de carga, demostrando que la resistencia interna del amperímetro perturba la medición (ej. alterando la corriente en $5.74\%$ en el rango de 50 mA versus $2.97\%$ en el de 100 mA), confirmando que el instrumento afecta al sistema original y que su precisión depende críticamente de una calibración adecuada para cada rango o encontrar correctamente los errores admisibles.